\section{Related Work}
\label{sec:related}

\subsection{Self-Stabilization in Distributed Systems}

Self-stabilization, introduced by Dijkstra~\cite{dijkstra_self_stab} and formalized
by Dolev~\cite{dolev_self_stab}, enables a distributed system to recover from any
transient fault within bounded steps without external intervention.
Classical results cover mutual exclusion, spanning trees, and clock synchronization under
crash faults.
Recent SSS proceedings extend the framework to adversarial and data-driven settings:
Chen, Lin, and Arora~\cite{sss_ml_routing} use ML predictions to improve the stabilization
time of a near-shortest-path routing protocol---a direct precedent for
learning-augmented self-stabilization.
Jobic et al.~\cite{sss_fl_leakage} analyze label leakage in federated learning using
cryptographic tools with formal security definitions.
Our work applies the stabilization framework to federated learning under Byzantine
(not crash) faults---qualitatively harder because Byzantine processes may deviate
arbitrarily rather than fail silently.

\subsection{Byzantine Fault Tolerance}

Byzantine fault tolerance~\cite{byzantine_generals} establishes that $f < n/3$ Byzantine
faults are tolerable in synchronous consensus; Castro and Liskov~\cite{bft_consensus}
extend this to partially synchronous systems.
The FLP impossibility~\cite{flp_impossibility} shows consensus is impossible with one
crash fault in asynchronous systems.
Recent SSS work examines Byzantine reliable broadcast~\cite{sss_byz_broadcast,sss_byz_counter}
and Byzantine-resilient blockchain sharding~\cite{sss_blockchain_sharding}.
Our protocol operates in the synchronous round model with $f < n/2$---weaker than the
$n/3$ BFT bound, enabled by our server architecture.

\subsection{Byzantine-Robust Federated Learning}

Krum~\cite{krum} selects gradients by nearest-neighbor distance but degrades under
Non-IID data~\cite{non_iid_challenges}.
Coordinate-wise robust statistics~\cite{byzantine_ml} reduce to $O(nd)$ computation
but are sensitive to gradient variance heterogeneity.
Bulyan~\cite{bulyan} iteratively filters outliers; FLTrust~\cite{crfl} uses a trusted
root dataset; FLAME~\cite{flame} targets backdoors via clustering; ByzShield~\cite{byzshield}
provides resilience via redundant training.
None specifies a formal recovery time or fault-containment invariant in the sense
of~\cite{dolev_self_stab}---motivating this work.

\subsection{Spectral Methods, RMT, and Blockchain FL}

The BBP phase transition~\cite{bbp_transition} characterizes when a rank-$f$ perturbation
creates eigenvalue outliers above the MP bulk edge---the mechanism underlying our detection
threshold.
Gradient covariance matrices of neural networks follow limiting spectral
distributions~\cite{rmt_gradients,rmt_neural,spectral_analysis}.
To our knowledge, we are the first to exploit the MP law~\cite{mp_law} for Byzantine
gradient identification with a formally scoped detection threshold connected to the BBP
phase transition.
Blockchain integration in FL~\cite{blockchain_fl,blockchain_smart} has focused on
auditability; we instead formalize the ledger as a self-stabilizing shared memory
substrate with Byzantine fault tolerance strictly stronger than classical shared registers.

\subsection{Positioning}

Our work differs from all of the above along three orthogonal axes.

\textbf{Self-stabilization vs.\ robustness.}
Byzantine-robust FL methods~\cite{krum,bulyan,crfl,flame,byzshield} guarantee convergence
from a \emph{known-good} initialization but offer no recovery guarantee from arbitrary
corrupted states.
Our protocol is the first to provide a formal recovery time bound
$T^* = O(\sigma^2\beta^2/\varepsilon^2 + \beta^2/\varepsilon)$ from any starting
configuration---satisfying the practical self-stabilization definition
of~\cite{dolev_self_stab} rather than merely stability under sustained attacks.

\textbf{Spectral mechanism vs.\ pairwise distance.}
Krum~\cite{krum}, Trimmed Mean~\cite{byzantine_ml}, Bulyan~\cite{bulyan}, and FLTrust~\cite{crfl}
all operate on pairwise or coordinate-wise distances between gradient vectors.
These metrics degrade at $\beta > 1/3$ because Byzantine clients can spread their updates
to match honest pairwise diversity.
Our spectral filter operates on the \emph{global eigenvalue distribution}, which remains
anomalous even when individual Byzantine updates appear locally plausible.
The detection threshold is derived analytically from the BBP phase transition~\cite{bbp_transition},
yielding the information-theoretic impossibility at $\sigma^2\beta^2 = 0.25$ as a
matching lower bound.

\textbf{Blockchain as stabilizing medium vs.\ auditability.}
Prior blockchain FL work~\cite{blockchain_fl,blockchain_smart} uses the ledger for
auditability and transparency.
We use it structurally: the write-once BFT consensus provides the Byzantine-fault-tolerant
analog of Dijkstra's shared registers~\cite{dijkstra_self_stab}, enabling the monotonic
progress argument (Lemma~\ref{lem:closure}) that classical shared memory cannot support
under Byzantine faults.
The closest SSS-community precedents---Byzantine reliable broadcast~\cite{sss_byz_broadcast,sss_byz_counter}
and blockchain sharding~\cite{sss_blockchain_sharding}---provide distributed safety
properties but do not address the self-stabilization of a learning protocol.
